\section{Row with Maximum Number of 1s}
\label{sec:bs-row-max-ones}

\noindent We now begin the third and final group of problems in this
chapter: \textbf{binary search on a 2D array}. The recurring theme is that
each \emph{row} (or column) of the matrix is independently sorted, which
lets us reuse the 1D binary search templates from
Sections~\ref{sec:bs-find-x}--\ref{sec:bs-peak-1d} \emph{per row or
column}, or combine row-sortedness and column-sortedness to search the
whole 2D index space in a single sweep.

\problemstatement{We are given an $n \times m$ binary matrix in which
\textbf{every row is sorted} (all $0$s appear before all $1$s within that
row). We must find the index of the row containing the \textbf{maximum
number of $1$s}.}

\begin{patternbox}
\emph{``Each row is sorted''} (even though the matrix as a whole is not
globally ordered in any simple way) is the keyword that tells us we can
binary search \emph{within each row independently}, rather than needing
any 2D-aware search technique. Whenever a 2D-array problem's sortedness
guarantee is row-by-row (or column-by-column) rather than global,
default first to ``run a 1D binary search on each row/column,'' since it
is almost always the simplest correct approach, even if not always the
asymptotically fastest one.
\end{patternbox}

\subsection*{Understanding the problem carefully}

Because each row is sorted with $0$s before $1$s, the number of $1$s in a
row equals $m - \textit{lowerBound}(\text{row}, 1)$, where
$\textit{lowerBound}(\text{row}, 1)$ is the first index in that row where
the value is $\ge 1$ (Section~\ref{sec:bs-lower-bound}). Computing this for
every row and tracking the maximum gives the answer.

\begin{ideabox}[Reuse, Don't Re-derive]
Apply Algorithm~\ref{sec:bs-lower-bound}'s \textsc{LowerBound} procedure
to each row independently (searching for the first $1$), convert each
result to a count of $1$s, and track the row index achieving the maximum
count.
\end{ideabox}

\subsection*{A note on the alternative $O(n+m)$ staircase approach}

A different, even faster technique exists for this specific problem:
starting from the top-right corner of the matrix, move left whenever the
current cell is $1$ (recording this row as the best-so-far and shrinking
the column boundary), and move down whenever the current cell is $0$. This
``staircase search'' exploits row-and-column-sortedness jointly and runs
in $O(n+m)$ total, faster than the $O(n \log m)$ of the per-row binary
search. We mention it here because it foreshadows the 2D search technique
used in Section~\ref{sec:bs-search-2d-matrix-2}, but the per-row binary
search remains the more directly transferable technique for this
chapter's purposes, since it requires no new derivation beyond what
Section~\ref{sec:bs-lower-bound} already established.

\subsection*{Algorithm (Per-Row Binary Search)}

\begin{enumerate}
  \item Initialize $\textit{bestRow} \gets -1$, $\textit{bestCount}
        \gets 0$.
  \item For each row $i$ from $0$ to $n-1$:
  \begin{enumerate}
    \item $\textit{idx} \gets \textit{lowerBound}(\text{row } i, 1)$.
    \item $\textit{count} \gets m - \textit{idx}$.
    \item If $\textit{count} > \textit{bestCount}$: update
          $\textit{bestCount} \gets \textit{count}$,
          $\textit{bestRow} \gets i$.
  \end{enumerate}
  \item Return \textit{bestRow}.
\end{enumerate}

\begin{algorithm}[H]
\caption{Row with Maximum Number of 1s}
\begin{algorithmic}[1]
\Procedure{RowWithMaxOnes}{\textit{matrix}}
  \State $\textit{bestRow} \gets -1$, $\textit{bestCount} \gets 0$
  \For{$i \gets 0$ to $n-1$}
    \State $\textit{idx} \gets \Call{LowerBound}{\textit{matrix}[i], 1}$
    \State $\textit{count} \gets m - \textit{idx}$
    \If{$\textit{count} > \textit{bestCount}$}
      \State $\textit{bestCount} \gets \textit{count}$, $\textit{bestRow} \gets i$
    \EndIf
  \EndFor
  \State \Return \textit{bestRow}
\EndProcedure
\end{algorithmic}
\end{algorithm}

\subsection*{Dry run}

\dryrun{Take the matrix
$\begin{bmatrix} 0&0&0 \\ 0&1&1 \\ 1&1&1 \\ 0&0&0 \end{bmatrix}$.}

\begin{itemize}
  \item Row $0$: $[0,0,0]$. Lower bound of $1$ is index $3$ (no $1$ found).
        Count $= 3-3=0$.
  \item Row $1$: $[0,1,1]$. Lower bound of $1$ is index $1$. Count
        $=3-1=2$. Best so far: row $1$, count $2$.
  \item Row $2$: $[1,1,1]$. Lower bound of $1$ is index $0$. Count
        $=3-0=3$. Best so far: row $2$, count $3$.
  \item Row $3$: $[0,0,0]$. Count $=0$.
\end{itemize}

The answer is row $2$, with $3$ ones.

\subsection*{C++ implementation}

\begin{lstlisting}
#include <bits/stdc++.h>
using namespace std;

int lowerBoundRow(vector<int>& row, int X) {
    int lo = 0, hi = row.size() - 1, ans = row.size();
    while (lo <= hi) {
        int mid = lo + (hi - lo) / 2;
        if (row[mid] >= X) { ans = mid; hi = mid - 1; }
        else lo = mid + 1;
    }
    return ans;
}

int rowWithMaxOnes(vector<vector<int>>& matrix) {
    int n = matrix.size(), m = matrix[0].size();
    int bestRow = -1, bestCount = 0;

    for (int i = 0; i < n; i++) {
        int idx = lowerBoundRow(matrix[i], 1);
        int count = m - idx;
        if (count > bestCount) {
            bestCount = count;
            bestRow = i;
        }
    }
    return bestRow;
}

int main() {
    vector<vector<int>> matrix = {
        {0,0,0},
        {0,1,1},
        {1,1,1},
        {0,0,0}
    };
    cout << "Row with max ones: " << rowWithMaxOnes(matrix) << endl;
    return 0;
}
\end{lstlisting}

\begin{complexitybox}
\textbf{Time Complexity.} For each of the $n$ rows, a binary search costs
$O(\log m)$, giving an overall time complexity of
\[
  O(n \log m).
\]
\textbf{Space Complexity.} $O(1)$ auxiliary space.
\end{complexitybox}

\begin{takeawaybox}
The simplest way to handle a 2D problem where sortedness is guaranteed
only \emph{per row} (or per column) is to apply a 1D binary search
template independently to each row, accepting an extra factor of $n$ (or
$m$) in the time complexity in exchange for needing no new derivation.
Faster joint row-and-column techniques exist for some such problems (the
staircase search above, and Section~\ref{sec:bs-search-2d-matrix-2}), but
they require the matrix to be sorted in \emph{both} directions
simultaneously.
\end{takeawaybox}
